
Neural theorem proving systems allocate substantial computational resources to proof search attempts. The research literature indicates that a significant fraction of proof attempts timeout without finding proofs. The LeanDojo benchmark paper (Yang et al., 2023) reports that 37\% of proof attempts timeout under standard resource budgets. The authors of the present work propose that more efficient resource allocation could be achieved through learned termination detection.

The key hypothesis is that LLM-based theorem provers produce confidence trajectories during proof search that encode geometric structure related to proof space familiarity. Specifically, the proposal is that entropy variance computed over softmax probability distributions can distinguish between searches likely to succeed versus those likely to timeout or diverge.

The research design includes three main components:

\begin{enumerate}
\item A confidence geometry framework that interprets entropy variance as a measure of manifold stability
\item A hybrid detector combining neural signals (confidence variance) with symbolic signals (state collisions)
\item An experimental protocol using extended-timeout experiments to generate ground truth labels
\end{enumerate}

The paper claims three contributions: establishing confidence geometry as a principle (validated by r=0.80 correlation), demonstrating a practical detector (F1=0.806 or F1=0.97 depending on which section of the paper is consulted), and revealing that simpler signal combinations outperform complex voting mechanisms.

\subsubsection{Critical Context on Empirical Validation}

Examination of the research artifacts reveals that the quantitative results reported throughout the paper are based on mock data generated by synthetic data generation functions. The file \texttt{h-m2/code/results/results_raw.csv} contains 100 entries with identifiers "mock_theorem_000" through "mock_theorem_099" and variance values generated according to predetermined statistical distributions specified in \texttt{h-m2/code/test_mock.py}. The mock data generation code creates variance values following normal distributions with means 0.09 (success) and 0.29 (timeout), which closely match the reported results.

A validation document dated 2026-04-20 (\texttt{h-m2/MOCK_FIX_SUMMARY.txt}) indicates awareness of this issue and claims that mock data was removed, but the results files examined still contain mock data entries. The experimental code includes placeholders for LeanDojo integration, but actual execution with real theorem proving appears not to have been completed.

This rewrite proceeds to describe the methodology as designed and the results as reported, but readers should understand that empirical claims have not been validated with real data.
